\section{אופרטורים טנזוריים ומשפט ויגנר-אקרט (\LR{Tensor Operators \& Wigner-Eckart Theorem})}

\subsection{מוטיבציה: חוקי טרנספורמציה של אופרטורים}
עד כה עסקנו בטרנספורמציה של וקטורים ומצבים קוונטיים תחת סיבוב. כעת נרחיב את הדיון לאופרטורים עצמם.
נזכיר כי וקטור $\vec{V}$ עובר טרנספורמציה תחת סיבוב $R$ באופן הבא:
\begin{equation}
    V'_i = \sum_j R_{ij} V_j
\end{equation}
במכניקת הקוונטים, אופרטור $O$ עובר טרנספורמציה תחת אופרטור הסיבוב היוניטרי $U(R)$ לפי:
\begin{equation}
    O' = U(R) O U^\dagger(R)
\end{equation}
המטרה שלנו היא להגדיר משפחה של אופרטורים, הנקראים \textbf{אופרטורים טנזוריים ספיריים} (\LR{Spherical Tensor Operators}), אשר מתנהגים תחת סיבובים בדיוק כמו המצבים העצמיים של התנע הזוויתי (כמו ההרמוניות הספיריות $Y_{lm}$).

\subsection{אופרטור טנזורי ספירי בלתי-פריק}
אנו מחפשים אופרטורים שניתן לסווג אותם לפי דרגה $\lambda$ (האנלוגית ל-$j$ או $l$) ורכיב $\mu$ (האנלוגי ל-$m$).

\begin{definitionbox}{הגדרה: אופרטור טנזורי ספירי}
סט של $2\lambda + 1$ אופרטורים, המסומנים $O_\mu^{(\lambda)}$ (כאשר $\mu = -\lambda, \dots, +\lambda$), נקרא \textbf{אופרטור טנזורי ספירי מדרגה $\lambda$} אם תחת סיבוב המוגדר על ידי מטריצות הסיבוב של ויגנר $D_{m'm}^{(j)}(R)$, הם עוברים טרנספורמציה לפי:
\begin{equation}
    U(R) O_\mu^{(\lambda)} U^\dagger(R) = \sum_{\mu'} D_{\mu'\mu}^{(\lambda)}(R) O_{\mu'}^{(\lambda)}
\end{equation}
כלומר, האופרטורים מתערבבים בינם לבין עצמם בדיוק כפי שמצבי תנע זוויתי $|j,m\rangle$ מתערבבים תחת סיבוב.
\end{definitionbox}

\subsubsection{הגדרה דיפרנציאלית (יחסי חילוף)}
לעיתים קרובות קשה לבדוק את חוק הטרנספורמציה המלא עבור סיבוב סופי. נוח יותר להשתמש בסיבובים אינפיניטסימליים, המובעים באמצעות יחסי חילוף עם יוצרי הסיבוב (אופרטורי התנע הזוויתי $\vec{J}$).

על ידי פיתוח אופרטור הסיבוב $U(\vec{n}, \theta) \approx 1 - \frac{i}{\hbar}\theta \vec{n}\cdot\vec{J}$ לזווית קטנה, ניתן להראות כי ההגדרה לעיל שקולה לסט יחסי החילוף הבאים:

\begin{theorembox}{יחסי חילוף של אופרטור טנזורי עם $\vec{J}$}
סט אופרטורים $O_\mu^{(\lambda)}$ הוא טנזור ספירי מדרגה $\lambda$ אם ורק אם מתקיימים היחסים הבאים:

\textbf{1. יחס חילוף עם $J_z$:}
\begin{equation}
    [J_z, O_\mu^{(\lambda)}] = \hbar \mu O_\mu^{(\lambda)}
\end{equation}

\textbf{2. יחס חילוף עם אופרטורי הסולם $J_\pm$:}
\begin{equation}
    [J_\pm, O_\mu^{(\lambda)}] = \hbar \sqrt{\lambda(\lambda+1) - \mu(\mu \pm 1)} O_{\mu \pm 1}^{(\lambda)}
\end{equation}
\end{theorembox}

\begin{notebox}{אנלוגיה למצבים קוונטיים}
שימו לב שהפעולה של הקומוטטור עם $J_\pm$ על האופרטור $O_\mu^{(\lambda)}$ זהה לחלוטין לפעולה של האופרטור $J_\pm$ על המצב העצמי $|\lambda, \mu\rangle$:
\[
J_\pm |\lambda, \mu\rangle = \hbar \sqrt{\lambda(\lambda+1) - \mu(\mu \pm 1)} |\lambda, \mu \pm 1\rangle
\]
זהו המפתח להבנת משפט ויגנר-אקרט: האופרטורים "נושאים" תנע זוויתי משלהם.
\end{notebox}

\subsection{משפט ויגנר-אקרט \LR{(Wigner-Eckart Theorem)}}
זהו אחד המשפטים החשובים ביותר בשימוש בסימטריות במכניקת הקוונטים. המשפט מאפשר לנו לחשב אלמנטי מטריצה של אופרטורים טנזוריים בין מצבים בעלי תנע זוויתי מוגדר.

\subsubsection{רעיון המכפלה הטנזורית}
כאשר אופרטור טנזורי $O_\mu^{(\lambda)}$ פועל על מצב קוונטי $|j, m\rangle$, המצב החדש שמתקבל:
\[
|\psi\rangle = O_\mu^{(\lambda)} |j, m\rangle
\]
מתנהג תחת סיבובים כמו מכפלה של שני תנעים זוויתיים: אחד בגודל $\lambda$ (מהאופרטור) ואחד בגודל $j$ (מהמצב).
לכן, ניתן לפרק את המצב הזה לצירוף ליניארי של מצבים בעלי תנע זוויתי כולל $J$ והיטל $M$, בעזרת מקדמי קלבש-גורדון (\LR{Clebsch-Gordan Coefficients}):
\[
O_\mu^{(\lambda)} |j, m\rangle = \sum_{J, M} C(\lambda, j, J; \mu, m, M) |\phi_{J, M}\rangle
\]

\subsubsection{ניסוח המשפט}
המשפט מפריד את אלמנט המטריצה לשני חלקים:
\begin{enumerate}
    \item \textbf{חלק גיאומטרי:} תלוי רק בסימטריה (במספרים הקוונטיים המגנטיים $m, \mu, m'$). חלק זה מוכל כולו במקדמי קלבש-גורדון.
    \item \textbf{חלק פיזיקלי/דינמי:} תלוי בטבע האופרטור והמצבים הרדיאליים, אך \textbf{אינו תלוי} במספרים המגנטיים. חלק זה נקרא "אלמנט המטריצה המצומצם" (\LR{Reduced Matrix Element}).
\end{enumerate}

\begin{resultbox}{משפט ויגנר-אקרט}
אלמנט המטריצה של אופרטור טנזורי $O_\mu^{(\lambda)}$ בין המצבים $|j, m\rangle$ ו-$\langle j', m'|$ נתון על ידי:
\begin{equation}
    \langle j', m' | O_\mu^{(\lambda)} | j, m \rangle = \langle j, \lambda; m, \mu | j', m' \rangle \frac{\langle j' || O^{(\lambda)} || j \rangle}{\sqrt{2j' + 1}}
\end{equation}
כאשר:
\begin{itemize}
    \item $\langle j, \lambda; m, \mu | j', m' \rangle$ הוא מקדם קלבש-גורדון (המחבר את התנעים $j$ ו-$\lambda$ לתנע כולל $j'$).
    \item $\langle j' || O^{(\lambda)} || j \rangle$ הוא אלמנט המטריצה המצומצם (המסומן לעיתים עם קווים כפולים או סוגריים מסולסלים). שימו לב: \textbf{הוא אינו תלוי ב-$m, \mu, m'$}.
\end{itemize}
\end{resultbox}

\begin{notebox}{הערה על סימונים}
ישנן קונבנציות שונות בספרות (למשל, הוספת פקטורים של $\sqrt{2j'+1}$ או מופעים של מינוס אחד). הנוסחה לעיל היא אחת הנפוצות (כפי שמופיעה אצל \LR{Sakurai} למשל), אך תמיד יש לוודא את ההגדרה הספציפית של אלמנט המטריצה המצומצם בספר בו משתמשים.
\end{notebox}

\subsection{יישום: כללי ברירה \LR{(Selection Rules)}}
הכוח הגדול של משפט ויגנר-אקרט הוא ביכולת לקבוע מתי אלמנט מטריצה \textbf{מתאפס זהותית} מבלי לחשב אינטגרלים מסובכים.

אלמנט המטריצה שונה מאפס רק אם מקדם הקלבש-גורדון שונה מאפס. מכאן נובעים שני תנאים הכרחיים (כללי ברירה):

\begin{enumerate}
    \item \textbf{שימור היטל התנע הזוויתי ($m$):}
    \[
    m' = m + \mu
    \]
    מקדם הקלבש-גורדון מתאפס אלא אם סכום ההיטלים נשמר.

    \item \textbf{אי-שוויון המשולש (תנע זוויתי כולל):}
    על מנת לחבר את תנע המצב ההתחלתי $j$ עם תנע האופרטור $\lambda$ ולקבל את המצב הסופי $j'$, חייב להתקיים:
    \[
    |j - \lambda| \le j' \le j + \lambda
    \]
\end{enumerate}

\subsubsection{דוגמה: אינטראקציה של אטום עם קרינה (קירוב דיפולי)}
\begin{examplebox}{מעברים אלקטרוניים באטום המימן}
נניח אטום מימן במצב יסוד $|n, l, m\rangle$ המקיים אינטראקציה עם פוטון. בקירוב של אורך גל ארוך ($e^{i\vec{k}\cdot\vec{r}} \approx 1 + i\vec{k}\cdot\vec{r}$), האינטראקציה הדומיננטית תלויה באופרטור המיקום $\vec{r}$.

הווקטור $\vec{r}$ הוא \textbf{אופרטור טנזורי מדרגה $\lambda=1$} (וקטור).
לכן, עבור מעבר ממצב $|l, m\rangle$ למצב $|l', m'\rangle$:

\textbf{1. כללי ברירה לתנע הזוויתי ($l$):}
לפי אי-שוויון המשולש עם $\lambda=1$:
\[
|l - 1| \le l' \le l + 1
\]
זה מאפשר $l' = l-1, l, l+1$.

\textbf{2. כללי ברירה לזוגיות (\LR{Parity}):}
האופרטור $\vec{r}$ הוא אי-זוגי תחת שיקוף ($\vec{r} \to -\vec{r}$).
פונקציות הגל של אטום המימן הן בעלות זוגיות $(-1)^l$.
כדי שהאינטגרל לא יתאפס, הזוגיות חייבת להשתנות (מכפלת זוגיות התחלה $\times$ אופרטור $\times$ סוף צריכה להיות זוגית):
\[
(-1)^{l'} \cdot (-1) \cdot (-1)^l = 1 \implies (-1)^{l' + l} = -1
\]
זה פוסל את המקרה $l' = l$.
\end{examplebox}

\begin{resultbox}{סיכום כללי ברירה למעבר דיפולי חשמלי}
\begin{itemize}
    \item $\Delta l = l' - l = \pm 1$ (המעבר $0 \to 0$ אסור).
    \item $\Delta m = m' - m = 0, \pm 1$ (תלוי בקיטוב הפוטון $\mu$).
\end{itemize}
מידע זה מתקבל ישירות מהסימטריה, ללא צורך בחישוב פונקציות הגל הרדיאליות!
\end{resultbox}

